\cotefile{157-5}
\scanpage{648}\mspage{36}
les foncteurs $u^{W}_{!}$. Notons déjà que \unsure{la théorie|le théorème} de
K.S. Brown, \inserted{sur les \unsure{instruments} de subdivision\textsuperscript{$*$},} donne essentiellement ceci :

% slanted
{\reversemarginpar\margin{\unsure{revus|revu} \deleted{\unsure{et}} \unsure{généralisés|généralisé} par mes soins}}

\uline{Proposition 3.2}. \deleted{Le prédérivateur défini par
$\mathrm{Ho}_{W}$, quand} Soit $W\inserted{\subset\mathrm{Fl}\,M}$ \deleted{est associé} un localiseur
associé à une structure à cofibrations
$(C,W)$ sur $M$. \deleted{\illegible{1}} \textcircled{\scriptsize 1$^\circ$} Sur la catégorie $\mathrm{Diag}_0$ des
ensembles ordonnés \uline{finis}, ce prédérivateur
satisfait aux axiomes Der 1 (décomposition),
Der 2 (localisation), Der 3 (objet final),
Der 4 bis (existence de $u^{W}_{!}$), Der 5 bis
(calcul des fibres de $u^{W}_{!}$), Der 6 bis
(ex. \unsure{strictes|suites} à gauche). \textcircled{\scriptsize 2$^\circ$} C'est encore vrai
sur la \inserted{\unsure{2-}}catégorie $\mathrm{Diag}_1$ \inserted{\unsure{de} Cat plus grande,} formée des ens.
ordonnés $I$ \deleted{\uline{noethériens}} tels que l'espace
top. associé à $I^\circ$ soit noethérien (\unsure{i.e.} tte
suite croissante d'objets de $I$, ou d'ouverts
de $I^\circ\deleted{\unsure{\}}}$, est stationnaire). \textcircled{\scriptsize 3$^\circ$} Enfin, si
\deleted{\unsure{$I$}} l'axiome transfini \inserted{Cof 4} est satisfait i.e.
si $C$ et $TC$ sont stables par composition
transfinie, alors les conclusions de $1^\circ$ restent
valables sur la \deleted{\unsure{cat.}} sous-catégorie \inserted{pleine} $\mathrm{Diag}_2$ de
Cat formée des ens ordonnés qui sont
noethériens sans plus.

\uline{Corollaire 3.3} Supposons que l'on ait Cof 3 sous
la forme plus forte d'une factorisation
% archive sheet number: 647
\scanpage{649}\mspage{37}
$f=pi$ \uline{fonctorielle} en $f$, ce qui est le cas
en particulier si la structure à cofibrations
est spéciale. Alors pour \unsure{tt|tte} $I$ dans
Cat, $\deleted{\unsure{M}}M(I)$ muni de $(W(I),C(I))$ est
une catégorie à cofibrations, laquelle \deleted{\unsure{ne peut}}
satisfait l'axiome transfini Cof 4 si
$(W,C)$ y satisfait. Donc on peut lui
appliquer les résultats de 3.2, qui
impliquent notamment que les axiomes
Der $1,2,3\,(?)$ sont valables sur Cat \unsure{tout entier|toute entière}, et que $u^{W}_{!}$ existe quand $u$
est de la forme $\mathrm{id}_{I}\times v$, avec $v:J\unsure{-|\to}K$
une flèche \unsure{dans} $\mathrm{Diag}_1$ (resp. dans $\mathrm{Diag}_2$, si
l'axiome transfini est satisfait).

3.4. La clef de la théorie est dans l'extension
de la structure à \unsure{cofibration|cofibrations} $(W,C)$, de $M$
à une catégorie $M(I)$, où $I$ est un ensemble
ordonné, en $(W(I),C_I)$. Les faits principaux
sont énoncés dans \inserted{$\underline{\overline{\mathrm{XVII}}}$, \unsure{aux|en}} \deleted{\unsure{les}} théorèmes p.~26--27, p.~50, p.~60.
C'est quand $I$ est \uline{noethérien}, et qu'on suppose
soit que $\mathrm{Esp}(I^\circ)$ est également noethérien, (p.ex.
$I$ fini), soit que $(W,C)$ satisfait l'axiome
transfini Cof 4, qu'on prouve que $(W(I),C_I)$
est également une structure à cofibrations sur $M(I)$,
% archive sheet number: 648

\scanpage{650}\mspage{38}
Un fait important pour la construction des
foncteurs $u^{W}_{!}$, pour \deleted{\unsure{$I$}} $u:I\longrightarrow I'$ une flèche
entre catégories ordonnées noethériennes, est
la relation
\[
 C_I\cap W(I)=(TC)_I,\tag*{(3.4.1)}
\]
d'où résulte que l'on a, pour $u$ comme dessus
\[
 u_!(C_I\cap W(I))\subset C_{I'}\cap\text{\uuline{$W(I')$}},\tag*{(3.4.2)}
\]
ingrédient essentiel de la construction du
foncteur $u^{W}_{!}$. Pour étendre cette construction
au cas d'une flèche $I$ quelconque dans Cat,
il se pose \deleted{\unsure{donc}} le problème d'étendre la
définition de $C_I$, et les formules (3.4.1) et
(3.4.2), au cas de cat. $\unsure{I|\ne}$ quelconque \unsure{dans|de} Cat.

Or on a prouvé

\uline{Proposition 3.5}. Si $\Phi\inserted{=\Psi^*}$ est une partie $Q$-\uline{close} à g.
de $\mathrm{Fl}\,M$, \deleted{\unsure{et}} et $I$ un ensemble ordonné
noethérien, $\Phi_I$ l'extension de $\Phi$ en une partie
de $\mathrm{Fl}\,M(I)$ définie par Anderson, alors on a
\[
 \Phi_I=\Psi(I)^*=v_!(\Phi(I_0)^{\sim},\tag*{\unsure{.|-}(3.5.1),}
\]
où $v:I_0\unsure{-|\to}I$ est le foncteur canonique de
la cat. discrète définie par $\mathrm{Ob}\,I$.

C'est (en des notations différentes) le
lemme 6.2.5. (p.~107) de $\underline{\overline{\mathrm{XVII}}}$, dont je ne
reprends pas ici la démonstration.

Ceci donne la clef d'une définition de $\Phi_I$
comme $\Phi\langle I\rangle$ (avec les notations des \S\S\,1,2),
% archive sheet number: 649
\scanpage{651}\mspage{\deleted{\unsure{38}}39}
pour $I$ quelconque dans Cat
\[
 \Phi_I\overset{\mathrm{def}}{=}\Phi\langle I\rangle\overset{\mathrm{def}}{=}\Psi(I)^*=v_!(\Phi(I_0))^{\sim},\tag*{(3.5.2)}
\]
où $\Psi=\Phi_*$, et où on suppose $\Phi$ Q-clos à g. i.e.
$\Phi=\Psi^*$. Dans le cas où on suppose
seulement $\Phi$ Q-saturé transfini, sans plus,
on aurait envie de poser
\[
 \Phi_I\overset{\mathrm{def}}{=}\text{Q-saturé transfini de }v_!(\Phi(I_0)),\tag*{(3.5.3)}
\]
et je présume que pour $I$ ordonné noethérien,
cela coïncide avec la définition d'Anderson.
Mais je vais me borner pour la suite au cas où
$C=\Phi$ est Q-clos à g.

Revenons alors au cas d'une structure à
cofibrations $(W,C)$ sur $M$, en supposant
$C$ et $TC=C\cap W$ Q-clos à gauche. Définissons
pour tout $I$ de Cat $C_I=C\langle I\rangle$ comme
en 3.5.2. On a alors, pour $u:I\unsure{-|\to}I'$ une
flèche de Cat
\[
 u_!(C\langle I\rangle)\subset C\langle I'\rangle,\tag*{(3.5.4)}
\]
mais il n'y a guère de raison d'espérer
avoir aussi (3.4.2), sans disposer de la
formule
\[
 C\langle I\rangle\cap W(I)\overset{?}{=}(TC)\langle I\rangle.\tag*{(3.5.5)}
\]
\margin{NB. On a l'inclusion $\supset$ ; car $v^*v_!TC(\unsure{I_0|I}\unsure{)}\subset TC(\unsure{I_1|I_0})$}
On a prouvé que cette formule est valable si $I$
est ordonné noethérien ($\underline{\overline{\mathrm{XVII}}}$ \inserted{th.~2.1} p.~26--27). Cela
% archive sheet number: 650
\scanpage{652}\mspage{40}
nous amène à la

\uline{Question 3.6} Soit $M$ une catégorie stable par \inserted{petites} $\lim$,
et soit $(\deleted{\unsure{C}}W,\deleted{\unsure{W}}C)$ une structure à cofibration dans
$M$ telle que $C$ et $TC=C\cap W$ soient Q-clos à gauche.
Soit $I$ un objet de Cat, a-t-on alors
\[
 C\langle I\rangle\cap W(I)\overset{?}{=}(TC)\langle I\rangle\;?\tag*{(3.6.1)}
\]

On a vu au \S\ précédent qu'il en est ainsi
si $(W,C)$ s'insère dans un triplet de
Quillen clos $(W,C,F=(TC)_*)$, \unsure{et quand on|et ayant en}
suppose de plus que les couples
\[
 C\langle I\rangle\longleftrightarrow TF(I),\qquad TC\langle I\rangle\longleftrightarrow F(I)\unsure{=}
\]
sont factorisants, auquel cas $(W(I),C\langle I\rangle,F(I))$
est une structure de Quillen close dans $M(I)$.
D'ailleurs, $(W,C,F)$ est une structure de Quillen
(néc. close) ssi les couples
$C\longleftrightarrow\underset{\inserted{\substack{\rotatebox{90}{$=$}\,\mathrm{def}\\ C_*}}}{\underbrace{TF}}$,
$TC\longleftrightarrow\underset{\inserted{\substack{\rotatebox{90}{$=$}\,\mathrm{def}\\ (TC)_*}}}{F}$
eux-mêmes sont factorisants\deleted{, et si de
plus $TF\subset W$ \unsure{(on aura alors}} et si de plus
$TF\subset W$ (ce qui est automatique si $W=W^{\unsure{\mathrm{G}|\leftrightarrow}}$,
car on aura en tous cas, par \deleted{\unsure{hypoth.}} Cof 4, $TF\subset W^{\mathrm{G}}$).

Mais peut-on s'en tirer avec moins, pour
pouvoir conclure à (3.6.1) ? Ne suffirait-il
pas de savoir que $TC\langle I\rangle\longleftrightarrow F(I)$ est
factorisant (ce qui n'a guère de chances
d'être vrai pour un $I$ $\unsure{\ne|I}$ général, si on n'a
déjà $TC\mathrel{\longleftarrow\mkern-9mu|\mkern-9mu\longrightarrow}F$, \unsure{factorisation|factorisant|factorisables} dans $M$
lui-même). C'est bien le cas :
% archive sheet number: 651

\scanpage{653}\mspage{41}
\uline{Proposition} 3.7 : \deleted{\unsure{Soit}} Sous les conditions générales de
(3.6) sur $M,W,C,I$, supposons que le
couple $TC\langle I\rangle\longleftrightarrow F(I)$ (où $F\overset{\mathrm{def}}{=}(TC)_*$)
soit factorisant, i.e. $TC\langle I\rangle\mathrel{\longleftarrow\mkern-9mu|\mkern-9mu\longrightarrow}F(I)$.
Alors la formule (3.6.1) est vraie.

En effet, si $f\in C\langle I\rangle\cap W(\unsure{J|I})$, on factorise $f$ en $pi$, avec
\diagram{\begin{tikzcd}[ampersand replacement=\&,column sep=2.4em,row sep=3.5em]
 X\arrow[rr,"{f\in C\langle I\rangle\cap W(J)}"]\arrow[dr,swap,"{i\in TC\langle I\rangle}"] \& \& Y\arrow[dl,dashed,bend right=12,"{\unsure{q|g}}"']\\
 \& X'\arrow[ur,bend right=12,swap,"{p\in F(I)}"] \&
\end{tikzcd}}
$i\in TC\langle I\rangle$, $p\in F(I)$. Comme
$f,i\in W(I)$, on aura
$p\in W(I)$ i.e. $p\in TF(I)$,
et à cause de $\underset{\inserted{\substack{\rotatebox{90}{$\in$}\\f}}}{C\langle I\rangle}\longleftrightarrow
\underset{\inserted{\substack{\rotatebox{90}{$\in$}\\p}}}{TF(I)}$,
il existe $q:Y\unsure{-|\to}X'$ telle que $pq=\mathrm{id}_Y$,
$qf=i$, donc $f$ est facteur direct de
$i\in TC\langle I\rangle$. Comme $TC\langle I\rangle$ est stable
par facteurs directs, on gagne.

\uline{Corollaire 3.8.} Soit $(\deleted{\unsure{C}}W,\deleted{\unsure{W}}C,F)$ une structure de
Quillen close dans $M$, et supposons que
$TC=C\cap W$ soit une partie Q-spéciale de $M$,
i.e. le couple $TC\mathrel{\longleftarrow\mkern-9mu|\mkern-9mu\longrightarrow}F$ un couple
spécial (i.e. $M$ est \uline{spéciale}, et $\exists$ petite
partie $TC_0$ Q-génératrice à g. pour $TC$).
Alors pour \unsure{tout|objet} $I$ \unsure{dans} Cat, on a
\[
 C\langle I\rangle\cap W(I)=TC\langle I\rangle.\tag*{(3.8.1)}
\]
% archive sheet number: 652

\scanpage{654}\mspage{42}
\deleted{\unsure{Donc|Dans}}
En apparence, on n'a pas eu à supposer
que $C$ soit elle aussi une partie Q-spéciale de
$\mathrm{Fl}\,M$. Mais comme on a besoin de la factorisation
pour \deleted{\illegible{1}} $C\longleftrightarrow TF$ (pour avoir au
moins une structure de cofibration sur $M$),
on n'a guère le moyen de l'établir sans
supposer $C$ spéciale. Sans compter que pour
pouvoir exploiter (3.8.1) pour la construction
de $u^{W}_{!}$, \unsure{on} \deleted{\unsure{faut}} \inserted{a} \deleted{supposer} \inserted{besoin} de savoir que
$(W(I),C\langle I\rangle)$ définit bien une structure
à cofibration sur $M(I)$, ce qui ne peut guère
se faire (il me semble) sans avoir la
factorisation pour $C\langle I\rangle\unsure{\longleftarrow|\longleftrightarrow}TF(I)$
également ; chose que je n'arrive à établir,
pour un $I$ général de Cat, que moyennant
l'hypothèse de spécialité sur $C$ également.

3.9. Pour en avoir le cœur net, il faudrait
quelques exemples concrets et illuminants, \unsure{du type suivant|des types suivants} :

a) Structure de Quillen \inserted{close} $(W,C,F)$ \inserted{sur $M$}, telle que pour
le prédérivateur associé, $u^{W}_{!}$ n'existe pas tjrs,
ou $u^{W}_{*}$ n'existe pas \unsure{toujrs|toujours}. Je pense
notamment au cas d'une flèche $u:I\unsure{-|\to}I'$
avec $I'=e$ (présumant que si dans ce cas,
$u^{W}_{!}$ resp. $u^{W}_{*}$ existe, alors il existe en tous les
cas). \unsure{Une|Un} \unsure{telle|tel} \deleted{\unsure{catastrophe|contrexemple}} serait encore plus
% archive sheet number: 653
\scanpage{655}\mspage{43}
instructif, bien sûr, si $M$ est stable par petites
$\varinjlim$, que ses objets soient accessibles etc, et
que $W$ soit \unsure{L-accessible}. (i.e. stable par
petites $\varinjlim$ filtrantes, et de plus accessible
\ldots). Le premier candidat pour $I$ auquel
je pense serait un $B_G$ ($G$ un groupe, p.ex.
$G=\pm1$), c'est en effet l'exemple aux antipodes
des ensembles ordonnés, lesquels, eux, ne
semblent pas faire \unsure{trop} de \unsure{problème|problèmes} (moyennant
l'hypothèse noethérienne \deleted{\unsure{tout}} au moins).

Les exemples suivants ne \deleted{\unsure{donnent}} sont pas
aussi convaincants que ça, hélas !

\uline{Exemple 3.10.} Soit $M$ une catégorie
quelconque, alors les couples $(\mathrm{Fl}\,M,\mathrm{Is}\,M)$
et $(\mathrm{Is}\,M,\mathrm{Fl}\,M)$ sont tous deux des
couples de Quillen
\[
 \mathrm{Fl}\,M\mathrel{\longleftarrow\mkern-9mu|\mkern-9mu\longrightarrow}\mathrm{Is}\,M
 \qquad \mathrm{Is}\,M\mathrel{\longleftarrow\mkern-9mu|\mkern-9mu\longrightarrow}\mathrm{Fl}\,M.\tag*{(3.10.1)}
\]
Il s'ensuit que si on pose
\[
 W=\mathrm{Is}\,M,\qquad \deleted{\unsure{E|F}}C=F=\mathrm{Fl}\,M,\tag*{(3.10.2)}
\]
$(W,C,F)$ est un \unsure{triple|triplet} de Quillen clos\textsuperscript{$*$}.
% slanted
{\reversemarginpar\margin{du moins si on admet que $M$ est stable par $\varinjlim$ et $\varprojlim$ \uline{finies}}}

D'ailleurs pour tout $I$ dans Cat, on a
$W(I)=\mathrm{Is}(M(I))$, donc \unsure{!}
\[
 \mathrm{Ho}_W(I)=M(I).\tag*{(3.10.3)}
\]

Donc la question de l'existence des foncteurs
$u^W_!,u^W_*$ se réduit à celle des foncteurs
$u_!,u_*$ habituels, adjoints \inserted{à g. et à dr.} de $u^*:M(I')\unsure{-|\to}M(I)$
% archive sheet number: 654

\scanpage{656}\mspage{44}
On sait bien que l'existence de ces foncteurs,
pour $u$ donné, fait appel à l'existence de
certaines $\varinjlim$ resp. $\varprojlim$ dans $M$ ; p.ex. si
$I'=e$, il est néc. et suffisant que les $\varinjlim$
(resp. $\varprojlim$) de type $I^\circ$ soient \unsure{repr.} dans $M$.
Mais avec les hypothèses faites sur $M$, il
n'y a pas de problème si $u:I\unsure{-|\to}I'$ est
un foncteur entre catégories \uline{finies}.

On va donner un exemple où il n'en
est plus ainsi.

\uline{Exemple 3.11} Soit $k$ un anneau
commutatif, et \deleted{\unsure{$\mathscr D$|$D$}} $M=K^b(k\text{-}\mathrm{Ab})$
la catégorie des complexes de \unsure{cochaines|cochaînes}
bornés, de $k$-modules, $W$ \unsure{les|la} \deleted{\unsure{constituée des}}
quasi-isomorphismes, donc
\[
 \mathrm{Ho}_W=\unsure{\mathrm{Der}|\mathrm{D}\deleted{\mathrm{er}}}^{b}(k)
\]
catégorie dérivée bornée de la catégorie
abélienne des $k$-modules. On sait,
en utilisant l'existence d'assez de
projectifs dans $k\text{-}\mathrm{Ab}$, que cette
structure est associée à une structure
de Quillen \inserted{close} $(W,C,F)$, décrite par Quillen
% archive sheet number: 655
\scanpage{657}\mspage{45}
dans Hom. Alg. -- à cela près qu'il \deleted{\unsure{en}} prend
$D^-(K)$ au lieu de $D^b(k)$. Peut-être,
pour vérifier les axiomes, doit-on supposer
que $k$ est de dim. coh. finie.

Prenons $I=B_G$, où $G$ est un
groupe (p.ex. $G=\{\pm1\}$), \deleted{Alors} et
$u:I\to e$. Alors $M(I)$ est la
catégorie des complexes bornés de
$k[G]$-modules. On voit que
$u^W_!$ resp. $u^W_*$ existe \unsure{sss|ssi} pour tt
tel complexe borné, ou simplement pour
un $k[G]$-module sans plus $M$,
le complexe $L\unsure{\Gamma|\mathscr T}^{G}(M)$ resp. $R\Gamma_G(M)$
\deleted{$\unsure{L\Gamma^{G}(M)}$} appartient à $D^b(k)$, ou ce
qui revient au $=$, sss on a
\[
 \begin{aligned}
 H_i(G,M)&=0 &&\text{pour }i\gg0\\
 \text{\deleted{\unsure{resp.}}}\quad H^i(G,M)&=0 &&\text{pour }i\gg0,
 \end{aligned}
\]
\unsure{tout} $k[\unsure{\ell|G|k}]$-module $M$. Ce n'est
pratiquement jamais le cas ; \deleted{pour} quand
on prend p.ex. $k=\mathbb Z$ et $G$ groupe fini,
cela n'est le cas que si $G=e$. L'exemple
le plus économique \deleted{\unsure{est}} \inserted{\unsure{paraît|serait}} $G=\pm1$, $k=\mathbb F_2$.
% archive sheet number: 656

\scanpage{658}\mspage{46}
Donc ici $u^W_!$ et $u^W_*$ n'existent pas, \unsure{pour}
$I\unsure{-|\to}e$ avec $I$ catégorie finie. Mais
aussi $M$ n'est pas stable par \inserted{petites} $\varinjlim$.

\uline{Exemple 3.12} Soit $M$ une catégorie munie
d'un localiseur $W$ (fortement saturé) tel
que $M-\underset{\inserted{\substack{\Vert\\\unsure{\mathscr H|\mathcal H}}}}{MW^{-1}}$ admette un adjoint
-- à gauche disons. On a donc une situation
de deux foncteurs adjoints
\diagram{\begin{minipage}{.92\linewidth}
\[
 \begin{tikzcd}[ampersand replacement=\&,column sep=6em]
 \mathscr H \arrow[r,shift left=1.0ex,"{v_!\ \text{pl. fid}}","{\scriptscriptstyle\vert}"{marking,pos=.55}] \&
 M \arrow[l,shift left=1.0ex,"{v^*\ \text{localisation}}"]
 \end{tikzcd}\,,\tag*{\unsure{(3.12 1)|(3.12.1)}}
\]
\end{minipage}}
et $W=(v^*)^{-1}(\mathrm{Iso}(\mathscr H))$. On trouve
que pour \unsure{tt|tte} $I$ dans Cat, \unsure{-- .|on a}
\[
 \mathrm H\deleted{\unsure{\mathrm O}}\mathrm o_W(I)\simeq\mathscr H(I)\overset{\mathrm{def}}{=}\mathrm{Hom}(I^\circ,\mathscr H).\tag*{\unsure{(3.17.2)|(3.12.2)}}
\]
On peut donc songer à faire un
exemple où dans $M$ les petites $\underset{\longleftarrow}{\varinjlim}$
existent, mais non dans $\mathscr H$ ; de sorte
que les $u^W_!$ ou $u^W_*$ n'existent pas
tjrs. Mais ça semble difficile à
construire, vu que dans la situation
envisagée, on voit facilement que les \inserted{petites}
$\varprojlim$ existent dans $\mathscr H$. Donc si dans $\mathscr H$
% archive sheet number: 657
\scanpage{659}\mspage{47}
il existe une petite famille cogénératrice\unsure{,|;}
(et pour ceci, il suffit qu'une telle
famille existe dans $M$), il y aura également
des \inserted{petites} $\varinjlim$, en vertu de SGA 4 I 8.12.7.
De même, elles existent si $\unsure{\mathscr H|\mathcal H}$ est une
sous-catégorie \inserted{pleine} de $M$ \uline{stable par} \inserted{\uline{petites}} $\varinjlim$ --
et il paraît pratiquement impossible
de faire un exemple (3.12.1) sans que
$\mathscr H$ ait cette propriété de stabilité.

Sans compter qu'il y a $=$ la perplexité
suivante :

\uline{Question 3.13} : Existe-t-il une
{\reversemarginpar\margin{!\hfill\rule[-2em]{.4pt}{4em}}}%
% vertical rule in the left margin beside the question
$\unsure{\mathcal U|\mathrm U}$-catégorie ($\mathcal U$ un univers donné)
où les \inserted{petites} $\varinjlim$ \deleted{quelconques} soient
représentables, sans qu'il en soit de
$=$ des petites $\varprojlim$ ? \uline{NB} Une
telle catégorie n'admet pas de petite
famille génératrice, a fortiori elle
n'est pas équivalente à une petite
catégorie.
% archive sheet number: 658

\scanpage{660}\mspage{48}
\uline{Question 3.14} Anderson prétend que pour
le prédérivateur défini par une catégorie
à cofibrations satisfaisant l'axiome Cof 4
transfini, les foncteurs $u^W_!$ existent.
Sa démonstration est \deleted{\unsure{tot}} \unsure{totalem\textsuperscript{t}|totalement}
\unsure{canulée}, néanmoins je n'arrive pas
à faire de contre-exemple. Le résultat
qu'il annonce serait-il vrai ? Je
ne connais pas $=$ d'exemple d'un
localiseur $W$, dans une catégorie $M$
\uline{avec \deleted{\unsure{f|fi}}petites} $\varinjlim$ \uline{et petites} $\varprojlim$,
tel que \deleted{\unsure{$u$}} $u^W_!$ ou $u^W_*$ n'existe\deleted{\unsure{nt}}
pas tjrs -- mais je suis persuadé
qu'il y a des exemples ($=$ avec $M$
une petite catégorie ?). J'en suis
moins convaincu si on suppose
$W$ stable par $\varinjlim$ filtrantes, et au
besoin, de plus, que les éléments de
$M$ sont accessibles (plus éventuellement
d'autres propriétés catégoriques anodines
sur $M$, telles l'existence d'une filtration
% archive sheet number: 659
\scanpage{661}
\mspage{49}
% archive sheet number: 660
cardinale etc).

\uline{Question 3.15} Cela me ramène à la question particulièrement \unsure{élusive} et irritante, de l’existence des petits ensembles $Q$-générateurs, dans une partie $\Phi$ de $\mathrm{Fl}\,M$, quand on suppose (disons) que $M$ est stable par petites $\varprojlim$, \inserted{que ses éléments sont accessibles} et \inserted{qu’il} admet une filtration cardinale.

Ainsi on a la question : une partie $L$-accessible de $\mathrm{Fl}\,M$, \inserted{\deleted{\unsure{q|p}}\unsure{disons -pré au besoin|disons -pour au besoin}} $Q$-saturée \deleted{\unsure{trans|traca}} à gauche (ou ce qui revient au $=$, contenant les iso, et stable par cochgt de base, par composition et par facteurs directs) admet-elle une \uline{petite} partie $\Phi_0$ $Q$-génératrice à g. ?

{\reversemarginpar\margin{!\hfill\rule[-3em]{.4pt}{6em}}}
D’autre part, je n’ai pas à présent construit un seul exemple d’un couple $\Phi\longleftrightarrow\Psi$, dans une catégorie $M$, qui ne soit factorisant ! Le cas où $M$ est spéciale, étant bien sûr le plus intéressant. Visiblement, dans tout
\scanpage{662}
\mspage{\textcircled{50}}
% archive sheet number: 661
ça ce sont les exemples qui manquent. La \unsure{clef} d’une compréhension semble bien être dans les questions de factorisation entre ensembles $\Phi\longleftrightarrow\Psi$\unsure{,|.} Il faudrait en faire une étude systématique, mais je recule devant la perspective d’une digression qui me prendrait peut-être encore des semaines !

\noindent\hrulefill\par\noindent\hrulefill

4. \inserted{\uline{Paratopos et}} \uline{L’opération} $M\boxtimes N$.

4.1. \inserted{\uline{Pseudo-topos.} \uline{\unsure{Pts. de première|Pts. du premier} Catégorie}} J’appelle \uline{pseudo-$\unsure{\mathcal U|\mathrm U}$-topos}, ou simplement \uline{pseudo-topos}, une $\mathcal U$-catégorie $M$ stable par $\mathcal U$-limites \deleted{\unsure{projectives}} \inserted{inductives}, et qui admet une \uline{petite} sous-catégorie $M_0$ \deleted{stable par} \uline{strictement génératrice} (i.e. génératrice par épimorphismes stricts).

Alors $M$ admet aussi des $\mathcal U$-limites projectives \inserted{(SGA 4 I 8.12.7)}. [\unsure{Mais} l’exemple de \unsure{(Gr)|(Gr.)}, cf loc cit 8.12.9, montre que $M^{\mathrm o}$ n’est pas forcément, pour \unsure{autant}, un pseudo-$\mathcal U$-topos, faute à $M$ d’admettre une petite \deleted{\unsure{famille}} sous-catégorie strictement \uline{cogénératrice}. Mais le th. 2.6. donne des conditions sur $M$ (\unsure{surtout} du type \unsure{exactitude})

